Antes de gerar o bloco do Capítulo 13, verifiquei manualmente a matemática por trás do pedido — e encontrei um problema real que preciso reportar antes de prosseguir com esse bloco específico.

## Achado: a fórmula "exata" de $d_s(\tau)$ não tem as assíntotas que o texto afirma

O pedido pede para eu "demonstrar" que $d_s(k) = 2 + \frac{2}{1+k/M_P}$ é o aproximante de Padé de $d_s(\tau)$ sob $\tau = 1/k^2$. Para isso, as duas expressões precisam compartilhar os mesmos limites assintóticos ($d_s \to 4$ no IR, $d_s \to 2$ no UV). Testei isso diretamente na fórmula existente (eq. `ds_tau`, linha 74, idêntica à eq. `analytical_ds_formula`, linha 417), usando expansões assintóticas padrão de $\operatorname{erfc}$ (Abramowitz–Stegun 7.1.23):

- **$\tau \to 0$** (expansão de Taylor de $\operatorname{erfc}(z) \approx 1 - 2z/\sqrt\pi$ em $z=\sqrt\tau/2$): $d_s(\tau) \to 4$ — não $2$, como o texto afirma nas linhas 76 e 421.
- **$\tau \to \infty$** (expansão de $\operatorname{erfcx}(z) \sim \frac{1}{z\sqrt\pi}(1-\frac{1}{2z^2}+\cdots)$): $d_s(\tau) \to 3$ — não $4$.

Ou seja, a fórmula fechada atual converge para **4 e 3**, não para **4 e 2** como reivindicado. Isso não é uma questão de terminologia — é um erro de cálculo no resultado central que o Cap. 12/13 chama de "primeira prova analítica rigorosa" do fluxo dimensional espectral. Escrever uma "demonstração exata" conectando essa fórmula a $2+2/(1+k/M_P)$ via Padé, como pedido, exigiria provar algo que é falso.

**Duas saídas honestas, sua escolha:**
1. Corrigir a fórmula de $d_s(\tau)$ (revisar a integral do kernel de calor / expoente) até que as assíntotas $4 \to 2$ realmente se verifiquem, e então eu derivo o Padé de verdade; ou
2. Reclassificar $d_s(k) = 2+\frac{2}{1+k/M_P}$ como um *ansatz fenomenológico independente* (motivado apenas por interpolar corretamente entre os limites qualitativos corretos, citando os resultados numéricos de CDT), sem alegar que é derivado exatamente de $d_s(\tau)$.

Não vou entregar o bloco CHAP13 até você decidir — não quero gerar um "QED" para uma conta que não fecha.

Os outros três pedidos são qualificações de terminologia/escopo que **são** sustentáveis com a literatura real, então seguem abaixo (mais um bônus: o Teorema 4.1 do Cap. 11, que precisa da mesma correção que o Teorema 5.1 do Cap. 12, já que ambos alegam equações de Einstein não-lineares completas a partir de um argumento que só é válido em primeira ordem).

%%% INICIO BLOCO CHAP12 JORDAN %%%
```latex
A persistent difficulty in canonical quantization is the presence of \emph{self-intersecting loops}: when $\gamma(s_1) = \gamma(s_2)$, the action of the Hamiltonian constraint operator $\widehat{\mathcal{H}}$ at the intersection vertex is subject to the regularization ambiguities intrinsic to Thiemann's construction --- namely the dependence of the operator's action on the \emph{valence} of the vertex created, the combinatorial \emph{Mandelstam branching} of the resulting Wilson-loop identities $W_{\gamma_1 \circ \gamma_2} = W_{\gamma_1} W_{\gamma_2}$, and the discretionary choice of \emph{node-splitting} prescription used to resolve the new edges.

\begin{theorem}[Jordan Embeddings and Regularization-Ambiguity-Free Loop States]
\label{thm:jordan_chronology}
Let $(\mathcal{M}^4, g)$ be a stably causal spacetime satisfying Hawking's Chronology Protection Conjecture.
\begin{enumerate}[label=(\roman*)]
    \item Any physical timelike worldline $\gamma(\tau)$ cannot self-intersect: $\gamma(\tau_1) \neq \gamma(\tau_2)$ for all $\tau_1 \neq \tau_2$.
    \item In any spacelike Cauchy surface $\Sigma \subset \mathcal{M}$, physical loop states generated by regularized holonomies must be \emph{Jordan curves} (injective topological circles $S^1 \hookrightarrow \Sigma$ without self-intersections).
    \item Restricting the quantum kinematic loop Hilbert space to Jordan curves removes the vertex-valence ambiguity, the Mandelstam-branching redundancy, and the node-splitting choice from the regularization of Thiemann's Hamiltonian constraint operator, so that on this restricted domain:
    \begin{equation}
        [\widehat{\mathcal{H}}[N], \widehat{\mathcal{H}}[M]] |\Psi_{\mathrm{Jordan}}\rangle = \widehat{\mathcal{H}}_a[\omega^a] |\Psi_{\mathrm{Jordan}}\rangle,
    \end{equation}
    i.e., the quantum hypersurface deformation algebra closes on Jordan-loop states without invoking a particular choice of vertex regularization scheme.
\end{enumerate}
\end{theorem}
\begin{proof}
Assertion (i) is a direct consequence of stable causality: a global time function $t \colon \mathcal{M} \to \R$ satisfies $\nabla_\mu t \dot{\gamma}^\mu > 0$. If $\gamma(\tau_1) = \gamma(\tau_2)$ with $\tau_1 < \tau_2$, the loop $\gamma|_{[\tau_1, \tau_2]}$ would be a Closed Timelike Curve (CTC), violating chronology.

By Section~1.2 of \cite{silvafilho2026noneuclidean}, the physical immersion problem and the Jordan embedding problem coincide identically. For (iii), Thiemann's regularization of $\widehat{\mathcal{H}}$ acts only at graph vertices; a self-intersection of $\gamma$ produces a vertex of valence $\ge 4$, whose action is regularization-scheme-dependent through the choice of loop assignment and edge-splitting used to resolve $\widehat{\mathcal{H}}$ at that node, and whose Wilson-loop content is constrained but not uniquely fixed by the Mandelstam identity $W_{\gamma_1 \circ \gamma_2} = W_{\gamma_1} W_{\gamma_2}$. For non-self-intersecting Jordan loops no such vertex is created, so none of these three regularization choices enter, and the constraint algebra closes independently of them.
\end{proof}
```
%%% FIM BLOCO CHAP12 JORDAN %%%

%%% INICIO BLOCO CHAP12 WDW %%%
```latex
\subsection{Minisuperspace Reduction of the Wheeler--DeWitt Constraint to the Closed-Loop Mode}
By Theorem~\ref{thm:jordan_chronology}, stable causality restricts physical quantum states to Jordan loops $S^1 \hookrightarrow \Sigma$. We consider the \emph{minisuperspace truncation} of the full Wheeler--DeWitt theory obtained by restricting attention to a single closed-loop mode $\gamma(s)$ and freezing all other kinematical degrees of freedom; this is the standard reduction used to extract a tractable one-dimensional model from the infinite-dimensional constraint, and it should not be read as a solution of the full multi-loop Hamiltonian constraint. Under the minimax extrinsic curvature bound $\|\II\|_{L^\infty} = \kappa^*$ (Theorem~\ref{thm:minimax_hamiltonian_regularization}), the functional Wheeler--DeWitt operator restricted to this single-mode sector reduces to a one-dimensional Sturm--Liouville eigenvalue equation:
\begin{equation}
    \label{eq:wdw_exact_sl}
    \left[ -\frac{\hbar^2}{\kappa^*} \frac{\dif^2}{\dif s^2} + V_{\mathrm{ADM}}(s) \right] \psi_\gamma(s) = 0.
\end{equation}
When the spatial curvature potential is linearized near the boundary of the embedding $V_{\mathrm{ADM}}(s) = V_0 + \kappa^* \sigma s$, equation \eqref{eq:wdw_exact_sl} has the exact analytical solution:
\begin{equation}
    \label{eq:wdw_airy_solution}
    \psi_\gamma(s) = c_1 \operatorname{Ai}\left( \left(\frac{(\kappa^*)^2 \sigma}{\hbar^2}\right)^{1/3} \left( s + \frac{V_0}{\kappa^* \sigma} \right) \right) + c_2 \operatorname{Bi}\left( \left(\frac{(\kappa^*)^2 \sigma}{\hbar^2}\right)^{1/3} \left( s + \frac{V_0}{\kappa^* \sigma} \right) \right),
\end{equation}
where $\operatorname{Ai}$ and $\operatorname{Bi}$ are the standard Airy functions. The boundary condition that $\psi_\gamma$ remains normalizable across the closed Jordan embedding enforces $c_2 = 0$, yielding discrete, quantized area eigenvalues \emph{within this single-loop minisuperspace sector}, free of the vertex-valence, Mandelstam-branching, and node-splitting ambiguities of Theorem~\ref{thm:jordan_chronology} (which do not arise here, since a Jordan loop contains no self-intersecting vertex).
```
%%% FIM BLOCO CHAP12 WDW %%%

%%% INICIO BLOCO CHAP12 HOLOGRAFIA %%%
```latex
\begin{theorem}[Wald Symplectic Equivalence: Linearized Einstein Equations, with Non-Linear Corrections from Relative-Entropy Positivity]
\label{thm:grand_einstein_emergence}
Let $|\Psi\rangle$ be a state in a continuous holographic tensor network (cMERA/cMPS), dual to a bulk geometry $g_{\mu\nu} = g_{\mu\nu}^{\mathrm{AdS}} + h_{\mu\nu}$ perturbatively close to pure AdS, with $h_{\mu\nu}$ treated to the order indicated below. Let $g^{\mathrm{QFI}}$ be the quantum Fisher information metric on the state manifold.
\begin{enumerate}[label=(\roman*)]
    \item \textbf{First order.} The first variation of the boundary first law of entanglement, $\delta S_A = \delta \langle H_A \rangle$, holding for all boundary spherical regions $A$, is equivalent via the Wald symplectic identity $\int_{\partial \Sigma} \chi = \int_\Sigma \dif \chi$ to the \emph{linearized} bulk Einstein equation
    \begin{equation}
        \delta\!\left( G_{\mu\nu} + \Lambda g_{\mu\nu} - 8\pi G_N \langle T_{\mu\nu}^{\mathrm{matter}} \rangle \right) = 0, \quad \forall (x, z) \in M,
    \end{equation}
    i.e.\ $h_{\mu\nu}$ satisfies the linearized Einstein equation to first order in the perturbation.
    \item \textbf{Second order and beyond.} The second variation of the modular relative entropy $S(\rho_A \| \sigma_A) \ge 0$ on boundary subregions $A$ reproduces the canonical energy of bulk metric perturbations. Its positivity, together with the monotonic relaxation of entanglement-cut surfaces under the level-set mean curvature flow of Theorem~\ref{thm:rt_mcf_convergence}, constrains the second-order and higher corrections to $h_{\mu\nu}$, extending the equation of motion order by order in the perturbative expansion.
\end{enumerate}
\end{theorem}
\begin{proof}
See Section~3 of \cite{silvafilho2026spacetime}. At first order, the Noether charge $(d-1)$-form on the extremal surface $\gamma_A$ integrates to $\delta \Area(\gamma_A)/(4G_N) = \delta S_A$, while on the boundary it integrates to the modular Hamiltonian $\delta\langle H_A\rangle$; the identity $\dif\chi = -2\xi^\mu \delta(G_{\mu\nu} + \Lambda g_{\mu\nu} - 8\pi G_N T_{\mu\nu})\mathbf{\epsilon}^\nu$ forces the linearized combination to vanish pointwise for all balls $A$, by the same sweeping argument as in Theorem~\ref{thm:einstein_emergence} of Chapter~11. This argument is intrinsically linear: it uses only the first variation of $S_A$, and therefore fixes $h_{\mu\nu}$ only to $\mathcal{O}(h)$. Extending the identification to $\mathcal{O}(h^2)$ and beyond requires the positivity of relative entropy $S(\rho_A\|\sigma_A) \ge 0$, whose second variation supplies the additional non-negative canonical-energy term fixing the non-linear completion of the equation of motion order by order, as in the entanglement-thermodynamics derivations of \cite{faulkner2014gravitation, lashkari2016canonical}. Full non-perturbative (all-orders) equivalence to the non-linear Einstein equations is not established by this argument alone.
\end{proof}
```
%%% FIM BLOCO CHAP12 HOLOGRAFIA %%%

**Nota**: as chaves `faulkner2014gravitation` (Faulkner, Guica, Hartman, Myers, Van Raamsdonk, *Gravitation from Entanglement in Holographic CFTs*, JHEP 2014) e `lashkari2016canonical` (Lashkari, Van Raamsdonk, *Canonical Energy is Quantum Fisher Information*, JHEP 2016) são placeholders — confirme/ajuste as chaves reais no seu `.bib`.

**Complemento necessário** — o Teorema 4.1 do Capítulo 11 (`thm:einstein_emergence`, linhas 208–260) tem exatamente o mesmo problema (conclui "full non-linear Einstein field equations" só a partir do first law linear) e precisa da mesma correção para manter consistência entre os capítulos:

%%% INICIO BLOCO CHAP11 EINSTEIN %%%
```latex
\begin{theorem}[Linearized Emergence of the Einstein Equations, with Non-Linear Corrections from Relative-Entropy Positivity]
\label{thm:einstein_emergence}
Assume:
\begin{enumerate}[label=(\alph*)]
    \item The holographic Ryu--Takayanagi relation holds for all boundary spherical subregions $A$:
    \begin{equation}
        S_A = \frac{\Area(\gamma_A)}{4 G_N},
    \end{equation}
    where $\gamma_A$ is the extremal bulk codimension-2 surface anchored at $\partial A$.
    \item The first law of entanglement entropy $\delta S_A = \delta \langle H_A \rangle$ holds for all balls $A$ of all radii $R > 0$ and centers $x_0 \in \R^{d-1, 1}$.
    \item The bulk matter fields satisfy the semiclassical stress-energy coupling $\langle T_{\mu\nu}^{\mathrm{matter}} \rangle$.
\end{enumerate}
Then the bulk metric perturbation $h_{\mu\nu} = g_{\mu\nu} - g_{\mu\nu}^{\mathrm{AdS}}$ satisfies, to first order in $h_{\mu\nu}$, the linearized Einstein field equations:
\begin{equation}
    \label{eq:einstein_emergent}
    \delta\!\left(G_{\mu\nu} + \Lambda g_{\mu\nu} - 8\pi G_N \langle T_{\mu\nu}^{\mathrm{matter}} \rangle\right) = 0, \quad \forall (x, z) \in M.
\end{equation}
The non-linear completion of \eqref{eq:einstein_emergent}, order by order in $h_{\mu\nu}$, follows from the positivity of the second-order modular relative entropy $S(\rho_A\|\sigma_A) \ge 0$ (Theorem~\ref{thm:grand_einstein_emergence} of Chapter~12).
\end{theorem}

\begin{proof}
Let $\chi$ denote the differential $(d-1)$-form defined in the bulk by the Wald symplectic form:
\begin{equation}
    \chi \coloneqq \delta \mathbf{Q}[\xi] - \xi \cdot \mathbf{\theta}(g, \delta g),
\end{equation}
where $\xi$ is the bulk Killing vector that generates the hyperbolic isometry preserving $\gamma_A$, with $\xi|_{\gamma_A} = 0$, and $\mathbf{Q}$ is the Noether charge $(d-1)$-form associated with diffeomorphism invariance.

By Stokes' theorem applied to the homology region $\Sigma_A$ between the boundary ball $A$ and the minimal surface $\gamma_A$ in the bulk:
\begin{equation}
    \int_{\partial \Sigma_A} \chi = \int_A \chi - \int_{\gamma_A} \chi = \int_{\Sigma_A} \dif \chi.
\end{equation}
On the boundary $z = 0$, the conformal Killing property of $\xi$ along with \eqref{eq:local_modular_ham} guarantees $\int_A \chi = \delta \langle H_A \rangle$. On the extremal surface $\gamma_A$, because $\xi|_{\gamma_A} = 0$ and its gradient gives the unit binormal, Wald's formula yields $\int_{\gamma_A} \chi = \delta \Area(\gamma_A)/(4 G_N) = \delta S_A$. Therefore the boundary first law $\delta S_A = \delta \langle H_A \rangle$ translates precisely into $\int_{\Sigma_A} \dif \chi = 0$. A direct calculation gives:
\begin{equation}
    \dif \chi = -2 \xi^\mu\, \delta\!\left( G_{\mu\nu} + \Lambda g_{\mu\nu} - 8\pi G_N \langle T_{\mu\nu}^{\mathrm{matter}} \rangle \right) \mathbf{\epsilon}^\nu,
\end{equation}
which is linear in the perturbation since $\delta S_A$ enters only at first order. Since this must vanish for all spherical regions $A$ of arbitrary center and radius, the fundamental lemma of the calculus of variations forces $\delta(G_{\mu\nu} + \Lambda g_{\mu\nu} - 8\pi G_N \langle T_{\mu\nu}^{\mathrm{matter}}\rangle) = 0$ pointwise, i.e.\ the linearized Einstein equation. This establishes the theorem at first order; the argument does not by itself constrain terms of $\mathcal{O}(h^2)$ or higher, which are fixed instead by the second-order relative entropy positivity argument of Chapter~12.
\end{proof}
```
%%% FIM BLOCO CHAP11 EINSTEIN %%%

Me diga como quer resolver o Capítulo 13 (opção 1: corrigir a fórmula de $d_s(\tau)$; opção 2: reclassificar $d_s(k)$ como ansatz independente) e eu gero o bloco final.
